% Candidate compact insertion for the Erdős 593 manuscript.
% This file is not included automatically.
% Integration requires adding a bibliography entry with key bahmanian2015.

\begin{remark}[Blocks and the earlier reduction]
\label{remark-komjath-block-reduction}
Komj\'ath proved that a finite triple system is obligatory if and only if all
of its $2$-connected components are obligatory
\citep{komjath2001}.  In the standard language of hypergraph connectivity,
these components are the blocks of the edge support; general block
decomposition and its block--cut tree are classical
\citep{bahmanian2015}.  Thus the content of
Theorem~\ref{theorem-canonical-atom-normal-form} is not the existence of a
block decomposition, but the exact identification of the blocks that can
occur in an obligatory triple system.
\end{remark}

\begin{corollary}[Block formulation of the classification]
\label{corollary-block-formulation}
Let $F$ be a finite triple system and delete its isolated points.  Then $F$ is
obligatory if and only if every block of $F^\circ$ is either one triple or
$J^+$ for a finite $2$-connected simple bipartite graph $J$.
\end{corollary}
\begin{proof}
If $F$ is obligatory, each block is obligatory by Komj\'ath's reduction and is
one-point indecomposable.  Corollary~\ref{corollary-minimal-generators-indecomposables}
therefore gives the displayed list.  Conversely every displayed block is
obligatory, and Komj\'ath's reduction gives obligatoriness of $F$.
\end{proof}

% Optional quantitative corollary.  This is the genuinely new part of the
% insertion; the partition-lattice algebra itself is classical.
\begin{corollary}[Exact one-point factorization-lattice spectrum]
\label{corollary-factorization-lattice-spectrum}
Let $F$ be reduced and obligatory, with $c\ge1$ connected components, shadow
order
\[
s=|V(F)|-|E(F)|,
\]
cycle rank
\[
\beta=2|E(F)|-|V(F)|+c,
\]
and $k$ canonical blocks.  Put
\[
N=k-c,
\qquad
q(r)=\left\lceil2\sqrt r\right\rceil.
\]
Let $\mathcal D(F)$ be the lattice, ordered by refinement, of decompositions
of the hyperedge set into connected supported pieces that meet pairwise in at
most one point and whose piece--shared-point incidence graph is a forest.

The possible ranks $N$ at fixed $(s,\beta,c)$ are exactly
\[
\begin{array}{ll}
\beta=0:&N=s-2c,\\[1mm]
\beta=1:&0\le N\le s-2c-2,\quad N\equiv s-2c\pmod2,\\[1mm]
\beta\ge2:&0\le N\le s-2c-q(\beta).
\end{array}
\tag{*}
\]
For every allowed $N$, the possible lattice isomorphism types are precisely
\[
\prod_{i=1}^t \Pi_{\lambda_i+1},
\qquad
\lambda=(\lambda_1,\ldots,\lambda_t)\vdash N,
\tag{**}
\]
where $\Pi_r$ is the partition lattice of an $r$-element set.  Different
partitions $\lambda$ give nonisomorphic lattices, and every lattice in
\((**)\) occurs.
\end{corollary}
\begin{proof}
The rank statement is the componentwise atom-count spectrum after subtracting
$c$.  If the shared points have multiplicities $\mu(p)$, the canonical block
forest gives
\[
N=\sum_p(\mu(p)-1).
\]
The bond lattice of the atom intersection block graph factors over its clique
blocks, one clique on $\mu(p)$ atoms for every shared point.  Hence
\[
\mathcal D(F)\cong\prod_p\Pi_{\mu(p)}.
\]
Writing $\lambda_i=\mu(p_i)-1$ gives \((**)\).  The exact shared-point
multiplicity theorem realises every partition $\lambda\vdash N$ without
changing the block list or the global parameters.  Finally the characteristic
polynomial of the product is
\[
\prod_i\prod_{j=1}^{\lambda_i}(x-j),
\]
whose root multiplicities recover $\lambda$, so distinct partitions give
nonisomorphic lattices.
\end{proof}
